% Chuzhditsa — manuscript for submission (XeLaTeX; requires the Chuzhditsa fonts in ../fonts/)
\documentclass[11pt]{article}
\usepackage[a4paper,margin=2.7cm]{geometry}
\usepackage{fontspec}
\usepackage{booktabs}
\usepackage{setspace}
\usepackage[hidelinks]{hyperref}

\setmainfont{Times New Roman}
\newfontfamily\ipafont{Arial Unicode MS}[Scale=MatchLowercase]
\newfontfamily\chu[
  Path = ../fonts/ ,
  Extension = .otf ,
  UprightFont = Chuzhditsa-Regular ,
  BoldFont = Chuzhditsa-Bold ,
  ItalicFont = Chuzhditsa-Italic ,
  BoldItalicFont = Chuzhditsa-BoldItalic ,
  Scale=MatchUppercase ]{Chuzhditsa-Regular}

% object language in the Chuzhditsa face; phonemic strings in an IPA-complete font
\newcommand{\cz}[1]{{\chu #1}}
\newcommand{\ph}[1]{{\ipafont/#1/}}
\newcommand{\phon}[1]{{\ipafont[#1]}}
\newcommand{\src}[1]{{\ipafont #1}}

\title{\textbf{Chuzhditsa: a featural loanword register for Cyrillic,\\ from Bulgarian to the Slavic superset}}
\author{Author name withheld for review\thanks{Fonts, build toolchain, and the twenty-language audit are available in the accompanying repository; the typeface used for all object-language examples in this manuscript is the system being described.}}
\date{}

\begin{document}
\maketitle

\begin{abstract}
\noindent Japanese distinguishes loanwords at a glance: katakana is a parallel register whose visual form alone announces foreignness, and whose adaptation rules map any foreign word to exactly one spelling. No Cyrillic orthography has an equivalent, and the result is systematic indeterminacy (Bulgarian \cz{Уошингтон}~$\sim$~\cz{Вашингтон}; Serbian and Russian counterparts abound). We present \emph{Chuzhditsa} (Bulg.\ чуждица, `foreignism'), a three-layer extension of Cyrillic that supplies (i) a featural grammar of diacritic operations deriving letters for non-Slavic phonemes, (ii) a prosodic layer for stress, tone, and length, and (iii) a register-marking typeface whose Glagolitic-derived forms perform the katakana function visually. The letter inventory is assembled conservatively: revivals from historical Bulgarian orthography (the yuses \cz{ѫ ѧ}) alongside the living Macedonian \cz{ѕ}, imports from the wider Cyrillic ecosystem (Kazakh \cz{қ ң}, Bashkir \cz{ҙ}, Tajik \cz{ӣ}), and rule-derived residue. We evaluate the system against the phonologies of twenty major languages, generalize it from its Bulgarian pilot to the full Slavic Cyrillic superset (with deterministic bridges from Latin-alphabet Slavic languages and an external validation on Albanian), and describe a parametric implementation as an OpenType typeface family in which the fusion and anchoring rules of the orthography are executable GSUB/GPOS code. We argue that the laryngeal series (\cz{х ҳ һ} plus the aspiration modifier \cz{◌ʰ}) illustrates a general principle for extending alphabets: articulatory gradients should map to visual gradients.
\end{abstract}

\section{Introduction}

Every literate community borrows words faster than its orthography can naturalize them. Japanese is unusual in having institutionalized the problem: katakana is a full parallel syllabary whose use \emph{is} the marking of foreignness, and whose adaptation conventions (\phon{naɪt}~$\rightarrow$~\src{ナイト}) are deterministic enough that two writers converge on one spelling. The Cyrillic-writing world has no analogous device. Loanwords enter Bulgarian, Russian, Serbian or Ukrainian through ad hoc transcription: sound-based and spelling-based renderings compete (\cz{ланч}~$\sim$~\cz{лънч} for \emph{lunch}), phonemes without Cyrillic letters collapse unpredictably (\ph{θ}~$\rightarrow$~т or с; \ph{w}~$\rightarrow$~в or у), and nothing in the written form signals that a word is a guest rather than a native.

This paper describes Chuzhditsa, a designed orthographic register that supplies the missing device for Cyrillic, piloted on Bulgarian and then generalized to the Slavic superset. The system makes three claims of general interest to writing-system design. First, that a loanword register can be \emph{featural} in Sampson's (1985) sense: rather than inventing letters ad libitum, a small set of diacritic operations — each meaning exactly one phonetic feature — derives the entire extension inventory, echoing the compositional logic of Hangul (Sampson 1985; Kim-Renaud 1997). Second, that such a register can be assembled almost entirely from attested material: dead letters of the target orthography's own history and living letters of related orthographies, with invention reserved for genuine residue. Third, that the register-marking function that Japanese distributes over a second character inventory can instead be carried by a second \emph{typeface register} — realized here as an OpenType family whose substitution (GSUB) and positioning (GPOS) rules are the orthography's fusion and attachment rules, executable rather than merely described.

\section{Background}

\subsection{Katakana as an adaptation layer}
Katakana's success rests on three properties (Irwin 2011; Stanlaw 2004): a fixed sound-mapping table applied to the source word's pronunciation, not its spelling (silent letters do not survive) — the orthographic reflex of loanword adaptation as phonological mapping (Kang 2011; Peperkamp \& Dupoux 2003); graphic distinctness sufficient to mark register at a glance; and extensibility (\src{ヴ} for \ph{v}, \src{ファ} for \ph{fa}) when the inventory runs out. Chuzhditsa adopts all three as requirements. Katakana also carries a social record that a designed analogue cannot disclaim: loanword marking has been read as a badge of incomplete nativeness, and katakana is deployed to mark foreigners' accented Japanese as well as foreign words (Irwin 2011; Robertson 2020). We therefore state the register's semantics narrowly: it asserts that a word's phonology lies off the native grid — nothing about its right to exist in the language, and nothing about its user. It marks words, never speakers, and a naturalized loan is expected to exit the register, exactly as \src{テンプラ} became \src{天ぷら}.

\subsection{The Cyrillic commonwealth}
Cyrillic serves some fifty languages (Comrie 1996), and the non-Slavic ones have already engineered letters for most sounds Slavic lacks: Kazakh \cz{қ}~\ph{q}, \cz{ң}~\ph{ŋ}, \cz{ғ}~\ph{ɣ}; Bashkir \cz{ҫ}~\ph{θ}, \cz{ҙ}~\ph{ð}; Tajik \cz{ӣ}~\ph{iː} and \cz{ҳ} (\ph{h} in Tajik, reassigned here to \ph{ħ} — the Caucasian digraph хӀ is the nearer precedent); Altai \cz{ӧ}~\ph{ø}, \cz{ӱ}~\ph{y}; Belarusian \cz{ў}~\ph{w}; Serbian \cz{џ}~\ph{dʒ}; the Caucasian palochka \cz{Ӏ}. These letters carry decades of typographic and pedagogic practice. A design that imports them inherits that practice for free — an option unavailable to katakana's designers and, we will argue, the single largest economy in the present system.

\subsection{The Glagolitic precedent}
Bulgaria is the locus of a script replacement: Glagolitic (863 CE), the first Slavic alphabet, was gradually displaced by Cyrillic (Preslav Literary School, \emph{c.}~893), which modelled most letters on Greek uncial and drew on Glagolitic mainly for the sounds Greek lacked (Schenker 1995; Cubberley 1996). Glagolitic then survived for a millennium as a special-register script (Croatian liturgical use into the twentieth century; the last Glagolitic missal was printed in 1905). Chuzhditsa knowingly reverses the vector: the visual DNA of the \emph{oldest} Slavic letterforms — round bowls, unicameral proportions — marks the \emph{newest} words, and the letters it revives (\cz{ѫ ѧ ѩ ѭ}) inherit the sound-slots of Glagolitic ancestors (the Cyrillic shapes themselves are reworkings).

\section{Design principles}

\begin{enumerate}
\item \textbf{Determinism.} One phoneme, one rendering; two transcribers converge. Loss is permitted but must be declared (\S\ref{sec:mergers}).
\item \textbf{Featural transparency.} Every diacritic operation means exactly one feature and applies wherever the feature appears.
\item \textbf{Pair symmetry.} Obstruent letters enter the system in voiceless/voiced pairs aligned with the native series. This principle outranks heritage: the historical theta-letter fita (\cz{ѳ}) was rejected because it has no voiced partner and derives from no rule; the Bashkir pair \cz{ҫ}/\cz{ҙ} (\ph{θ}/\ph{ð}) does its work systematically.
\item \textbf{Precedent first.} Revive native dead letters, then import living Cyrillic letters, and only then derive new combinations; never invent an unattested base shape.
\item \textbf{Pan-Cyrillic legibility.} Every base skeleton follows the reader's existing Cyrillic, and every extension must be guessable at first sight by any Cyrillic reader from its base letter plus a transparent mark — a design criterion the orthography-development literature calls transferability (Cahill \& Rice 2014).
\item \textbf{Register by typeface.} Foreignness is marked by letterform register, not by punctuation or a second inventory: text set in the Chuzhditsa face is thereby marked, exactly as katakana text is — a register distinction of the kind the biscriptality literature treats as script-internal diglossia (Sebba 2007; Bun\v{c}i\'c, Lippert \& Rabus 2016).
\end{enumerate}

\section{The featural grammar}

Table~\ref{tab:grammar} gives the full operation set. The retraction hook generalizes a precedent internal to the Slavic base inventory itself — \cz{щ} \emph{is} \cz{ш} with a hook — and the glide breve generalizes \cz{и}~$\rightarrow$~\cz{й}. Where Unicode supplies precomposed hooked letters (\cz{қ ң ҳ ҫ ҙ}) they are used; elsewhere the same hook is written with combining U+0322 (\cz{т̢ д̢ р̢}).

\begin{table}[ht]\centering\small
\begin{tabular}{lll}
\toprule
Mark & Operation & Derivations \\
\midrule
\cz{◌̆} breve & vowel $\rightarrow$ glide & \cz{и→й} \ph{j} · \cz{у→ў} \ph{w} \\
\cz{◌̢} hook below & place shifted off the native grid & \cz{қ} \ph{q} · \cz{ҳ} \ph{ħ} · \cz{ң} \ph{ŋ} · \cz{ҫ} \ph{θ} · \cz{ҙ} \ph{ð} · \cz{т̢} \ph{ʈ} · \cz{д̢} \ph{ɖ} · \cz{р̢} \ph{ɽ} \\
\cz{◌̵} crossbar & lenition & \cz{г→ғ} \ph{ɣ} \\
\cz{◌̈} double dot & vowel fronting & \cz{ӓ} \ph{æ} · \cz{ӧ} \ph{ø œ} · \cz{ӱ} \ph{y ʏ} \\
\cz{◌̨} yus tail & nasalization & \cz{ѫ} \ph{ɔ̃ õ} · \cz{ѧ} \ph{ɛ̃} · \cz{а̨} \ph{ɑ̃} · \cz{ӧ̨} \ph{œ̃} \\
\cz{◌ʰ ◌ʱ} raised \cz{һ} & aspiration; breathy voice & \cz{пʰ тʰ кʰ чʰ} · \cz{бʱ дʱ гʱ џʱ} \\
\cz{◌̄} macron & vowel length & \cz{ӣ ӯ} (precomposed) · \cz{а̄ е̄ о̄} \\
\cz{Ӏ} palochka & laryngeal--pharyngeal action & \cz{тӀ кӀ цӀ} (ejectives) · bare \cz{Ӏ} = \ph{ʕ} \\
\cz{ь} & palatalization & \cz{нь} \ph{ɲ} · \cz{ль} \ph{ʎ} · \cz{хь} \ph{ç} · \cz{сь} \ph{ɕ} \\
\bottomrule
\end{tabular}
\caption{The featural operations. Each mark has one meaning; precomposed Unicode letters are used where they exist.}
\label{tab:grammar}
\end{table}

\subsection{The laryngeal gradient}
The three back fricatives that plague Cyrillic transcription of Arabic, English and German — \ph{x}, \ph{ħ}, \ph{h} — receive a designed gradient: articulatory weakening maps to graphic lightening. Native \cz{х} (\ph{x}$\sim$\ph{χ}) is two full crossing strokes, the densest friction; pharyngeal \cz{ҳ} (\ph{ħ}) is \cz{х} pulled deeper by the retraction hook; glottal \cz{һ} (\ph{h}$\sim$\ph{ɦ}) is a single open arc, literally the most open glyph of the three; and aspiration \cz{◌ʰ}, a feature rather than a segment, is \cz{һ} reduced to a superscript. Arabic exhibits the full triple within ordinary vocabulary: \cz{һила̄л} \ph{hilaːl}, \cz{Муҳаммад} \ph{muħammad}, \cz{Ха̄лид} \ph{xaːlid}. We suggest the mapping \emph{articulatory gradient $\rightarrow$ visual gradient} as a reusable principle for alphabet extension.

\subsection{The Indic stress test}
Hindi--Urdu and Bengali cross a four-way laryngeal contrast with a dental/retroflex place contrast. The grammar covers the entire grid with zero new letters: \cz{т}~\ph{t̪}, \cz{тʰ}~\ph{t̪ʰ}, \cz{д}~\ph{d̪}, \cz{дʱ}~\ph{d̪ʱ}; retroflex \cz{т̢ т̢ʰ д̢ д̢ʱ}, plus \cz{н̢}~\ph{ɳ}, \cz{р̢}~\ph{ɽ}, \cz{р̢ʱ}~\ph{ɽʱ}. A single word such as \cz{т̢ʰӣк} (\emph{ṭhīk}, \ph{ʈʰiːk}) composes three features on one base — hook, aspiration, length — which is precisely the featural claim made concrete.

\section{The ligature stratum}

Cyrillic has always ligated: \cz{ю} is a historic ligature of і+оу (the у element later dropped), Serbian \cz{љ њ} were fused from л/н+ь by Vuk Karad\v{z}i\'c in 1818, building on Mrkalj's digraphs, and Old Bulgarian possessed iotated nasals. Chuzhditsa's ligature stratum is therefore conservative, and it is implemented as executable GSUB rules rather than prose: (i)~\cz{нь ль} fuse to њ/љ-forms automatically (\cz{ниньо}, \cz{фамилья} set with fused palatals); (ii)~the iotated yuses \cz{ѩ}~\ph{jɛ̃} and \cz{ѭ}~\ph{jɔ̃} are revived as both directly typeable letters and automatic fusions of \cz{ьѧ}/\cz{ьѫ} (French \emph{bien}~$\rightarrow$~\cz{бѩ}; Polish \emph{pięć}~$\rightarrow$~\cz{пѩч}); (iii)~ejective digraphs \cz{тӀ кӀ пӀ цӀ чӀ} fuse with the palochka \emph{raised to the upper half}, echoing the apostrophe of scholarly romanization (\phon{tʼ kʼ}) — a fused full-height palochka was rejected in testing because \cz{тӀ} became confusable with \cz{п}. The governing rule: only sequences denoting a single phoneme, with Slavic or Old Cyrillic precedent, may fuse; mechanical soft pairs (\cz{сь хь}) remain visibly two-part.

\section{The prosodic stratum}

Stress takes the grave of Bulgarian lexicography (\cz{ба̀бушкъ}); Mandarin tones 2--4 take acute, caron, grave, with tone 1 unmarked — pinyin minus the macron, which Chuzhditsa reserves for length. The unmarked first tone resolves the collision by construction: \cz{ма̄} can only mean a long vowel. Three costs are declared rather than hidden. Tone and stress marks are homoglyphic across source languages, resolved only by knowing the lexical source (katakana shares exactly this property). Diacritic stacking degrades: Vietnamese \emph{phở} wants quality, length and tone at once (\cz{фъ̄̌}), and the running-text rule drops length first (\cz{фъ̌}). And four marks cannot represent six-tone systems without folding (Vietnamese \emph{hỏi/ngã} merge). The design verdict generalizes katakana's own: prosody is a dictionary layer, not an orthographic one.

\section{From Bulgarian to the Slavic superset}
\label{sec:panslavic}

The pilot treated Bulgarian's thirty letters as the base. Nothing in the featural grammar is Bulgarian-specific, however, and the generalization is administrative rather than structural: the base becomes each language's own Cyrillic alphabet, and the extension and prosodic strata apply unchanged. Concretely, the typeface's character set was widened to the Slavic Cyrillic superset — Russian \cz{э ё ы}, Ukrainian \cz{і ї є ґ}, Belarusian \cz{ў}, Serbian \cz{ј љ њ ђ ћ џ}, Macedonian \cz{ѕ ѓ ќ} — so that any Slavic Cyrillic text sets natively. One design conflict surfaced and is instructive: the pilot's round е-form was, in fact, the shape of Ukrainian \cz{є}; pan-Slavic coverage forced е to an angular form so that the е/є contrast survives. Extending a script for one language can silently occupy a sibling's letterform space; superset-first design catches this early.

For each language the system's marginal contribution differs, and several languages turn out to need the extension layer for \emph{native} material, not only loans: Slovak \emph{ä} is \cz{ӓ} and Slovak \emph{h} is \ph{ɦ}~$\rightarrow$~\cz{һ}; Slovenian's native schwa (\emph{polglasnik}) is exactly \cz{ъ} (\cz{пъс} \ph{pəs} `dog'); Serbian's pitch accents are the prosodic layer's use case; Polish \emph{ą ę} are the yuses meeting their own descendants (\cz{мѫка}, \cz{пѩчь}). Latin-alphabet Slavic languages connect through a deterministic bridge: Gaj's digraphs map one-to-one onto Serbian Cyrillic letters (\emph{lj nj dž}~$\rightarrow$~\cz{љ њ џ}), and the residue is featural — Czech \emph{ř} \ph{r̝}, the notorious singleton, is written \cz{р̌} with the caron it already wears; Slovak long syllabic liquids \emph{ĺ ŕ} take the length macron (\cz{вр̄ба} \ph{vr̩ːba}).

The generalization also has ancestors to acknowledge: the determinism principle is Vuk Karad\v{z}i\'c's \emph{пиши као што говориш} avant la lettre, and Serbian orthography contributed \cz{џ љ њ ј} — and the very idea of fused palatals — before this project existed; Ukrainian orthography constrained the design as a peer, not a client (the є letterform conflict above was resolved in Ukrainian's favor). The Bulgarian name records the pilot, not a hierarchy.

As an external stress test we applied the system to Albanian, a non-Slavic Balkan language: every Albanian phoneme lands on an existing letter (\emph{ë}~\ph{ə}~$\rightarrow$~\cz{ъ}, \emph{th/dh}~$\rightarrow$~\cz{ҫ}/\cz{ҙ}, \emph{y}~$\rightarrow$~\cz{ӱ}, \emph{xh}~$\rightarrow$~\cz{џ}, \emph{x}~\ph{dz}~$\rightarrow$~\cz{ѕ}, \emph{q/gj}~$\rightarrow$~\cz{кь}/\cz{гь}): \cz{Шкьипъри́а}, \cz{Тира́нъ}. A system built for Slavic phonology plus the featural operations appears to cover the Balkan sprachbund without amendment. The practical beneficiary is not Albanian itself, which is securely Latin-written, but Cyrillic-language text citing Albanian: the \emph{ë} that Macedonian and Serbian renderings currently drop is exactly \cz{ъ}.

\section{Typeface implementation}

The Chuzhditsa family (Regular, Bold, Italic, Bold Italic; TTF and CFF/OTF) is generated by a $\sim$300-line parametric builder: every glyph is a set of monoline strokes — lines, arcs, and the full circles that give the face its Glagolitic character — expanded to outlines by a union of per-segment rectangles and per-vertex discs, with weight and slant as parameters. Capitals stand at 700 units, lowercase at an x-height of 500 with true two-case architecture (bowl-and-stem \cz{а}, ascending \cz{б}, descending \cz{р}). Every combining mark receives computed GPOS mark-to-base anchors (the macron of \cz{а̄} attaches 251 units lower than that of \cz{А̄}; the retroflex hook of \cz{т̢} lands on the letter's stem), and the ligature stratum is a GSUB \texttt{liga} feature — the orthography's rules are the font's code, verified by HarfBuzz shaping tests in the build toolchain — a concrete instance of Sproat's (2000) treatment of orthography as a computable mapping. The manuscript the reader is holding sets all object-language material in this typeface.

\section{Evaluation: the twenty-language audit}

The system was audited against the phoneme inventories of the twenty most-spoken languages (clicks excluded by scope). Table~\ref{tab:examples} samples the audit; the full 126-item set accompanies the repository. Declared mergers\label{sec:mergers} are: all rhotics~$\rightarrow$~\cz{р}; \ph{ʂ}/\ph{ʃ}~$\rightarrow$~\cz{ш}; \ph{ɸ}/\ph{f}~$\rightarrow$~\cz{ф}; \ph{ɦ}/\ph{h}~$\rightarrow$~\cz{һ}; implosives~$\rightarrow$~\cz{б д}; \ph{q}/\ph{ɢ}~$\rightarrow$~\cz{қ}; Korean tense consonants as geminates; English input referenced to a rhotic accent (GA); \ph{ʔ} dropped outside Arabic material.

\begin{table}[ht]\centering\small
\begin{tabular}{llll}
\toprule
Source & & IPA & Chuzhditsa \\
\midrule
English & \emph{thanks} & \phon{θæŋks} & \cz{ҫӓңкс} \\
English & \emph{weekend} & \phon{ˈwiːkɛnd} & \cz{ўӣкенд} \\
Mandarin & \src{北京} \emph{Běijīng} & \phon{pèɪ.tɕíŋ} & \cz{Пе̌йчиң} \\
Hindi & \src{ठीक} \emph{ṭhīk} & \phon{ʈʰiːk} & \cz{т̢ʰӣк} \\
Arabic & محمد \emph{Muḥammad} & \phon{muˈħammad} & \cz{Муҳаммад} \\
French & \emph{croissant} & \phon{kʁwasɑ̃} & \cz{крўаса̨} \\
French & \emph{bien} & \phon{bjɛ̃} & \cz{бѩ} \\
Polish & \emph{Wrocław} & \phon{ˈvrɔt͡swaf} & \cz{Вроцўаф} \\
Japanese & \src{東京} \emph{Tōkyō} & \phon{toːkʲoː} & \cz{То̄кьо̄} \\
Korean & \src{한글} \emph{hangul} & \phon{hanɡɯl} & \cz{һангыл} \\
Georgian & \src{წყალი} \emph{c\textquotesingle qali} & \phon{t͡sʼqʼali} & \cz{цӀқӀали} \\
Albanian & \emph{Shqipëria} & \phon{ʃcipəˈɾia} & \cz{Шкьипъри́а} \\
\bottomrule
\end{tabular}
\caption{Sampled audit items; the Chuzhditsa column is set in the typeface under discussion, with live GSUB fusion and GPOS anchoring.}
\label{tab:examples}
\end{table}

\section{Adoption: lessons from planned and organic systems}

A designed register lives or dies socially, not technically; the sociolinguistics of script choice behaves like the sociolinguistics of language choice (Fishman 1977; Unseth 2005). Four historical cases bound the space Chuzhditsa enters, and they are worth reading as an engineer reads failure reports.

\textbf{The planned failures.} Volap\"uk collapsed within a decade, less on linguistic defects than on its inventor's proprietary control of the standard (Garv\'ia 2015). Esperanto — maximally learnable, energetically organized — plateaued as a voluntary subculture: it extended nobody's existing practice and attached to no prior community, so adopting it meant \emph{joining a new identity} rather than exercising an old one (Forster 1982; Garv\'ia 2015). The Deseret alphabet is the sharpest lesson, because it controls for institutional power: commissioned by Brigham Young, taught in Utah schools, backed by church funds, it was abandoned within roughly two decades — sponsorship without a felt need or an identity payoff bought nothing (Alder, Goodfellow \& Watt 1984). Learnability and patronage, the two levers a designer can pull, decided none of these cases.

\textbf{The designed success.} Hangul is the one engineered script that won, and its path is instructive precisely because the design was not what won. Completed in 1443 and promulgated in 1446 (Ledyard 1998), it was opposed at once by the literati — the collective Ch'oe Malli memorial of 1444 — and spent four and a half centuries outside prestige writing, surviving in low-stakes registers: women's correspondence, vernacular fiction, private devotion (Lee \& Ramsey 2011). What revalued it was not its featural elegance but nineteenth-century nationalism, culminating in the 1894 Kabo edict that made the national script the base of official documents (King 1998). The mandate arrived 448 years after the design; the niche did the surviving.

\textbf{The organic success.} Katakana never asked anyone to change anything. It began as monks' abbreviated annotation glosses on Chinese texts in the early Heian period, drifted into a division of labor with hiragana over a millennium, and was codified only in 1900 by the implementing regulations of the Elementary School Order — with the loanword-register specialization consolidating later still (Seeley 1991; Twine 1991). Codification ratified practice; it did not create it.

The pattern reduces to four conditions, against which Chuzhditsa can be scored honestly. (i) \emph{Additive, not substitutive}: the winners changed no existing spelling; Chuzhditsa alters nothing in any native orthography by construction. (ii) \emph{Identity attachment}: Hangul survived on Korean-ness, not on design merit; Chuzhditsa's revivals — the yuses, the Glagolitic letterforms — are the analogous move available to a designed system, heritage rather than invention. (iii) \emph{A survivable low-stakes niche before any mandate}: Hangul had the vernacular genres, katakana the annotation margin; the dictionary pronunciation field, the language classroom, and the subtitle are proposed here as exactly that — the system's vernacular centuries, compressed. (iv) \emph{Codification last}: every mandate that worked ratified existing practice, which is why this specification petitions no orthographic authority. One condition cannot be engineered around: on the planned/organic axis Chuzhditsa sits with Esperanto and Deseret, not with katakana. Assembly from attested material is a mitigation — every letter arrives with a history already attached — but it is not an exemption, and the Deseret case fixes an upper bound on what design quality plus sponsorship can buy without a community that recognizes itself in the script.

\section{Limitations}

The retraction hook is polysemous by direction — uvular on \cz{к}, pharyngeal on \cz{х}, interdental on \cz{с з}, retroflex on \cz{т д р}. No attested language contrasts two of its values on one base letter, so the ambiguity is theoretical, but it is a real weakening of the one-mark-one-meaning claim. \ph{ɕ} keeps two spellings (\cz{сь} under contrast, \cz{ш} by entrenchment) — the one point where usage beat system. Six-tone prosodies fold lossily into four marks. Click phonology is future work, as is a mark-to-mark (\texttt{mkmk}) positioning pass for the rare double-stacked vowel.

\section{Conclusion}

Chuzhditsa demonstrates that the katakana function — deterministic adaptation plus visible register — can be retrofitted onto an alphabet without inventing a parallel inventory: a featural grammar over attested Cyrillic material covers twenty phonologies, one historical script supplies the register aesthetics, and the whole specification compiles to a font whose OpenType rules \emph{are} the orthography. The generalization from Bulgarian to the Slavic superset cost one letterform (е) and no grammar changes, which we take as evidence that the featural layer, not the letter list, is the right unit of design.

\begin{thebibliography}{99}
\bibitem{alder} Alder, D.~D., P.~J.~Goodfellow \& R.~G.~Watt. 1984. Creating a new alphabet for Zion: The origin of the Deseret alphabet. \emph{Utah Historical Quarterly} 52(3). 275--286.
\bibitem{buncic} Bun\v{c}i\'c, D., S.~L.~Lippert \& A.~Rabus (eds.). 2016. \emph{Biscriptality: A Sociolinguistic Typology}. Heidelberg: Universit\"atsverlag Winter.
\bibitem{cahill} Cahill, M. \& K.~Rice (eds.). 2014. \emph{Developing Orthographies for Unwritten Languages}. Dallas, TX: SIL International.
\bibitem{comrie} Comrie, B. \& G.~G.~Corbett (eds.). 1993. \emph{The Slavonic Languages}. London: Routledge.
\bibitem{comrie96} Comrie, B. 1996. Adaptations of the Cyrillic alphabet. In P.~Daniels \& W.~Bright (eds.), \emph{The World's Writing Systems}, 700--726. New York: OUP.
\bibitem{coulmas} Coulmas, F. 2003. \emph{Writing Systems: An Introduction to their Linguistic Analysis}. Cambridge: CUP.
\bibitem{cubberley} Cubberley, P. 1996. The Slavic alphabets. In P.~Daniels \& W.~Bright (eds.), \emph{The World's Writing Systems}, 346--355. New York: OUP.
\bibitem{danchev} Данчев, А. 1982. \emph{Българска транскрипция на английски имена}. 2.~изд. София: Народна просвета.
\bibitem{danielsX} Fishman, J.~A. (ed.). 1977. \emph{Advances in the Creation and Revision of Writing Systems}. The Hague: Mouton.
\bibitem{forster} Forster, P.~G. 1982. \emph{The Esperanto Movement}. The Hague: Mouton.
\bibitem{garvia} Garv\'ia, R. 2015. \emph{Esperanto and Its Rivals: The Struggle for an International Language}. Philadelphia: University of Pennsylvania Press.
\bibitem{daniels} Daniels, P. \& W.~Bright (eds.). 1996. \emph{The World's Writing Systems}. New York: OUP.
\bibitem{irwin} Irwin, M. 2011. \emph{Loanwords in Japanese}. Amsterdam/Philadelphia: John Benjamins.
\bibitem{kimrenaud} Kim-Renaud, Y.-K. (ed.). 1997. \emph{The Korean Alphabet: Its History and Structure}. Honolulu: University of Hawai`i Press.
\bibitem{king98} King, R. 1998. Nationalism and language reform in Korea: The questione della lingua in precolonial Korea. In H.~I.~Pai \& T.~R.~Tangherlini (eds.), \emph{Nationalism and the Construction of Korean Identity}, 33--72. Berkeley: Institute of East Asian Studies.
\bibitem{ledyard} Ledyard, G.~K. 1998. \emph{The Korean Language Reform of 1446}. Seoul: Sin'gu Munhwasa.
\bibitem{leeramsey} Lee, K.-M. \& S.~R.~Ramsey. 2011. \emph{A History of the Korean Language}. Cambridge: CUP.
\bibitem{kang} Kang, Y. 2011. Loanword phonology. In M.~van Oostendorp, C.~J.~Ewen, E.~Hume \& K.~Rice (eds.), \emph{The Blackwell Companion to Phonology}, vol.~4, 2258--2282. Malden, MA: Wiley-Blackwell.
\bibitem{peperkamp} Peperkamp, S. \& E.~Dupoux. 2003. Reinterpreting loanword adaptations: The role of perception. \emph{Proceedings of the 15th International Congress of Phonetic Sciences}, 367--370.
\bibitem{robertson} Robertson, W.~C. 2020. \emph{Scripting Japan: Orthography, Variation, and the Creation of Meaning in Written Japanese}. London: Routledge.
\bibitem{sampson} Sampson, G. 1985. \emph{Writing Systems: A Linguistic Introduction}. Stanford: Stanford University Press.
\bibitem{sebba} Sebba, M. 2007. \emph{Spelling and Society: The Culture and Politics of Orthography around the World}. Cambridge: CUP.
\bibitem{schenker} Schenker, A.~M. 1995. \emph{The Dawn of Slavic: An Introduction to Slavic Philology}. New Haven: Yale University Press.
\bibitem{seeley} Seeley, C. 1991. \emph{A History of Writing in Japan}. Leiden: Brill.
\bibitem{sproat} Sproat, R. 2000. \emph{A Computational Theory of Writing Systems}. Cambridge: CUP.
\bibitem{stanlaw} Stanlaw, J. 2004. \emph{Japanese English: Language and Culture Contact}. Hong Kong: Hong Kong University Press.
\bibitem{twine} Twine, N. 1991. \emph{Language and the Modern State: The Reform of Written Japanese}. London: Routledge.
\bibitem{unseth} Unseth, P. 2005. Sociolinguistic parallels between choosing scripts and languages. \emph{Written Language \& Literacy} 8(1). 19--42.
\bibitem{unicode} The Unicode Consortium. 2024. \emph{The Unicode Standard}, Version 16.0.0. South San Francisco, CA: The Unicode Consortium.
\end{thebibliography}

\end{document}
